\chapter*{Preface}
A theorem is learned twice: once when it is proved, and again when it is needed somewhere nobody said it would be. The
five volumes of the first part of \emph{Order and Chance} prove their theorems chapter by chapter, and each chapter's
exercises live next door to the tools they need. This book is the second lesson. Its three hundred problems are
organised by meeting points, not by tomes: roots of unity, finite fields, random integers, Fourier analysis, symmetry,
asymptotics. A problem about dice is solved by factoring polynomials; a question about Fibonacci numbers is settled in a
field with $p^2$ elements; the determinant of a matrix of greatest common divisors turns out to be a product of Euler's
$\varphi$; Beukers's double integrals prove that $\zeta(3)$ is irrational.

The statements never say which chapter applies. Finding the right tool is half of each problem, and the most important
half: it is the skill that research needs and that no single chapter can teach. When the search stalls, the hints at the
end of each chapter give a nudge, sometimes by naming the tome where the tool is proved. Every solution then starts with
\emph{How to find it}, which tells how one might think of the approach, including the false starts worth knowing, and
continues with a complete \emph{Write-up}. A short line of \emph{Tools used} names the tomes and chapters that were
needed, and the map at the end of each chapter shows all the crossings at once.

Problems are graded \textbf{$\star\star$} for a good afternoon and \textbf{$\star\star\star$} for a gem that may take a
weekend; a few \textbf{$\star$} problems open each chapter, and the problems marked \emph{Journey} are long walks in many
steps that end at a famous theorem. Some of the problems go on to research questions, which the remarks \emph{Where it
leads} describe, with references; results quoted there for their interest are labelled as such and never used in a proof.
Everything that is used is proved here or in a named tome.


The pleasure of a problem book is the moment when two distant ideas click together. We have tried to put as many of those
moments as possible between these covers.
